\begin{center}
\textsc{Gauss} to \textsc{Gerling}.\quad Göttingen, 5~September 1823.
\end{center}

\par\noindent\hspace*{\parindent}\mbox{.\ .\ .\ .\ .\ .}\quad I definitely want you to put the heliotrope [on the
Inselsberg] exactly in line with the theodolite site and to pitch the tent
as symmetrically as possible. If you could give the side facing the
Brocken a dark colour, so much the better. I
\textls[120]{hope} that the heliotrope there will far outshine the tent. But I
did have a case in 1822 when the \textls[120]{white} tent in Scharnhorst made my
observations in Garssen unusable for an entire day. It should be noted,
though, that first the tent turned its sunlit southern side towards me,
and second the heliotrope light was very weak. Neither condition
applies here. Still, as a precaution I would like the heliotrope
\textls[120]{not to be placed off-centre}, since otherwise the tent could contribute
slightly to distorting the angle. I suspect that if, for example, the
tent is \(\frac{1}{40}\) as bright as the heliotrope and the latter is far enough away for
the separation to be \(4^{\prime\prime}\), the angle will be distorted by \(0\mathord{\ooalign{\hfil\(\vphantom{0}^{\prime\prime}\)\hfil\cr\hfil\({,}\)\hfil\cr}}1\).\\*
\mbox{.\ .\ .\ .\ .\ .}\par
